\chapter{インフレーション、構造、そして初期条件の喪失}

\section{オメガと平坦性の証拠}

大きな\(\Omega\)についての学生の質問は、議論をフリードマン・ロバートソン・ウォーカー方程式へと戻す。放射優勢のかなり後である現在の時代については、膨張則は真空エネルギーの項と物質の項を用いて書かれ、放射は今では非常に小さいので落とされる：
\begin{equation}
H^2=\left(\frac{\dot a}{a}\right)^2=\frac{8\pi G}{3}\left(\rho_0+\rho_m\right)-\frac{k}{a^2}.
\end{equation}
真空の項は薄まらない。物質の項は薄まる：
\begin{equation}
\rho_m=\frac{M}{a^3}.
\end{equation}
\(\Omega\)は、膨張方程式のうち異なる成分が担う割合である。

\(H^2\)で割る。すると規格化された項の和は1になる。正則化された記法では、
\begin{equation}
\begin{aligned}
\Omega_0&=\frac{8\pi G\rho_0}{3H^2},\\
\Omega_m&=\frac{8\pi G\rho_m}{3H^2},\\
\Omega_k&=-\frac{k}{a^2H^2},
\end{aligned}
\end{equation}
したがって
\begin{equation}
\Omega_0+\Omega_m+\Omega_k=1.
\end{equation}
\editorialnote{この区間は音声も黒板も粗く、話者は真空の寄与をいくつかの名前で呼び分けている。ここでの記法は、\(H^2\)で割った後の議論を正則化するものである。}
空間が平坦ならば\(k=0\)であり、したがって
\begin{equation}
\Omega_0+\Omega_m=1.
\end{equation}
物質と真空の和が1より大きければ、曲率は正である。1より小さければ、曲率は負である。

観測的な論理はそれから単純である。通常の物質と暗黒物質を測る。ハッブル膨張を測る。非常に大きなスケールで曲率を幾何学的に測る。さらに、宇宙が加速していることにも気づく。真空の項の符号はそれ自体が\(H\)であることではなく、\(\dot a/a\)が定数へ向かう傾向であり、それはスケール因子が指数関数的な枝
\begin{equation}
a(t)\propto e^{Ht}.
\end{equation}
に近づくことを意味する。いくつかの測定は互いに照合し合い、宇宙は数パーセント以内で非常に平坦に近いと言う。

\begin{classroomqa}
\lecturetimestamp{00:10:50}
\audiencequestion{物質には暗黒物質が含まれるのか。}
\lecturerresponse{はい。現在の議論では、物質の項には暗黒物質が含まれる。放射を加えることもできるが、今日その割合は無視できるほど小さい。}
\end{classroomqa}

\begin{classroomqa}
\lecturetimestamp{00:11:11}
\audiencequestion{ハッブル定数以外に暗黒エネルギーの尺度はあるのか。}
\lecturerresponse{要点は\(H\)の値そのものではなく、膨張が指数関数的な法則へと向かう仕方である。物質の測定および平坦性の幾何学的測定と合わせて、それによって真空の項を推論できる。}
\end{classroomqa}

\section{平坦性から構造へ}

なぜ宇宙はそれほど平坦なのか。その提案がインフレーションである。平坦とは、何十億光年、実に何百億光年のスケールで、非常に大きく、非常に滑らかで、非常に一様であることを意味する。指数関数的膨張は、ある領域をそのように見えるまで引き伸ばすことができる。

しかし、良いことでもやり過ぎることがある。完全に滑らかな宇宙は完全に滑らかなままでいる。重力は厳密な対称性からむらを作り出さない。今構造が存在するならば、構造はもっと早い時期の小さな非一様性から始まらなければならなかった。現在の宇宙を過去へたどると、透明になった時点での非一様性はおよそ\(10^5\)分の1になる：
\begin{equation}
\frac{\delta\rho}{\rho}\sim 10^{-5}.
\end{equation}
宇宙は極めて滑らかであったが、完全にそうであったわけではない。

そこから重力は、空気で満たされた部屋の中の通常の力がしないことをする。小さな過密度は物質を自分自身へと引き寄せる。近くの領域は枯渇するので、低密度はさらに低密度になり、過密度はさらに過密度になる。重力はコントラストを消すのではなく強める。

自然な疑問は、なぜこれがもっと早く起こらなかったのかである。答えは圧力である。もっと早い時期には、放射圧が大きかった。
\begin{equation}
p=w\rho,
\end{equation}
を用いると、よく知られた場合は
\begin{equation}
w_{\rm rad}=\frac13,\qquad
w_m=0,\qquad
w_{\rm vac}=-1.
\end{equation}
である。正の圧力は一様化する。それは揺らぎを押し出し、媒質を滑らかにする。大きなスケールでは無圧力として扱われる物質は、そのような支えを提供しない。圧力が存在する間は、重力は抑えられる。圧力が無視できるようになると、重力が勝つ。

暗黒物質は出来事の順序において決定的である。最初に崩壊する塊は暗黒物質ハローである。次に通常のバリオン物質がそれらのハローへ落ち込み、バリオンは散逸できるので中心により近く集まる。暗黒物質は同じようには散逸しないので、ハロー全体に広がったままである。その収縮の始まりは透明になる時点よりわずかに早く、時間でおよそ10倍であり、それがマイクロ波背景放射がその始まりを直接示さない理由である。

短い補足がその後の化学史を確定する。ヘリウムより重い元素のほとんどすべては、初期宇宙の宇宙論ではなく恒星進化の産物である。初期宇宙は本質的に水素、ヘリウム、そしていくつかの軽い原子核の微量の痕跡を作る。

空で観測されるのは質量密度そのものではなく、マイクロ波背景放射の温度異方性である。密度揺らぎがそれに寄与する。速度揺らぎも寄与する。過密度領域はわずかに低温に見える傾向がある。光子がより深い重力井戸を登り抜ける際に余分なエネルギーを失うからである。それに加えて、ガスの運動によるドップラーシフトが重なる。完全な数値解析は複雑である。頑健に言えることはもっと単純である。宇宙が透明になった時点で、エネルギー密度は約\(10^5\)分の1だけ変動していた。

\begin{classroomqa}
\lecturetimestamp{00:20:29}
\audiencequestion{なぜ重力はもっと早くに塊を成長させなかったのか。}
\lecturerresponse{それ以前は宇宙が放射優勢であり、放射が媒体を均質化するのに十分な正の圧力を持っていたからである。放射が物質に比べて重要でなくなったときに初めて、重力が塊の成長を支配し始めた。}
\end{classroomqa}

\begin{classroomqa}
\lecturetimestamp{00:24:35}
\audiencequestion{収縮には散逸機構が必要ではないのか。}
\lecturerresponse{通常のバリオン物質が中心へ落ち着くには確かに散逸が必要である。しかし暗黒物質は、その種の衝突的散逸なしにハローを作ることから始められる。膨張そのものに結びついた宇宙論的な種類の減衰もあり、それは後で重要になる。}
\end{classroomqa}

\begin{classroomqa}
\lecturetimestamp{00:30:34}
\audiencequestion{速度揺らぎもあり、後の銀河が角運動量を獲得できたのか。}
\lecturerresponse{はい。原始的なパターンには密度揺らぎと速度揺らぎが含まれていた。空に見えるのは、密度変動とガスの運動の両方によって生じた温度パターンである。}
\end{classroomqa}

\section{インフラトンとプラトー}

ここで導入される機構はスカラー場\(\phi\)である。そのエネルギー密度には三つの部分がある。
\begin{equation}
\rho_\phi=\frac12\dot\phi^2+\frac12(\nabla\phi)^2+V(\phi).
\end{equation}
時間変化が寄与する。空間変化が寄与する。スカラー場は真のポテンシャルエネルギー\(V(\phi)\)も担うことができる。その部分は異例ではない。ヒッグス場が身近な例である。ここで重要なのはポテンシャルの形である。

ポテンシャルはある範囲にわたってかなり平坦で、わずかに下向きの傾きを持ち、その後、臨界値で底が抜けなければならない。それはちょうどゼロまで落ちるのではない。非常に小さい何かまで落ちる。プランク単位では、現在の真空エネルギーは程度
\begin{equation}
10^{-123}.
\end{equation}
である。プラトーの上ではこれはこれよりずっと大きかったが、それでもプランクスケールよりはるかに低かった。

\begin{figure}[t]
\centering
\includegraphics[width=0.78\textwidth]{lecture_08_figure_03.png}
\caption{広いプラトー、急な落下、低い極小を持つ定性的なインフラトンポテンシャル。}
\lectureframe{00:43:23}
\end{figure}

このプロットは空間内の場所を記述するものではない。場の値を記述するものである。宇宙がプラトーの高いところから始まったと言うのは、非常に大きな領域にわたって場の値がそこにあったと言うにすぎない。最初に\(\phi\)が場所ごとにいくらか変動していても重要ではない。インフレーションがそれらの変動を途方もなく引き伸ばすからである。十分な膨張の後には、どの局所的な領域もほとんど一様に見える。\editorialnote{黒板のフレームでは軸ラベルは判読できないが、口頭の議論から、縦軸がポテンシャルエネルギー\(V(\phi)\)であり、横軸が場の値\(\phi\)であることは明らかである。}

場がポテンシャルの広く緩やかに傾斜した部分にとどまる限り、ポテンシャルエネルギーは真空エネルギーのように振る舞う。もし\(\phi\)が完全に凍結していたなら、それはちょうど真空エネルギーのように振る舞うだろう。省略された数因子を別にすれば、講義で述べられた関係は単に
\begin{equation}
H\propto \sqrt{V(\phi)},
\end{equation}
であり、それによりほぼ指数関数的な法則
\begin{equation}
a(t)\propto e^{Ht}.
\end{equation}
が得られる。現在観測できるほとんどすべてのことは、場が縁を越えた後に起こる。落下は、プラトー上の大きな真空様エネルギーが放射、通常物質、暗黒物質、そして今日見られるわずかな残存真空エネルギーに変換される場所である。

\begin{classroomqa}
\lecturetimestamp{00:44:07}
\audiencequestion{\(\phi\)は何の上の場であり、アインシュタイン方程式はまだ使われているのか。}
\lecturerresponse{はい。\(\phi\)は時空上の場であり、同じフリードマン方程式がまだ使われている。この領域では、右辺のエネルギー密度は主に\(V(\phi)\)によって供給される。ただし場が十分ゆっくり変化する場合である。}
\end{classroomqa}

\begin{classroomqa}
\lecturetimestamp{00:47:04}
\audiencequestion{真空様のインフラトンエネルギーはどのようにして物質と放射に変わりうるのか。}
\lecturerresponse{場が底まで転がって振動すると、それらの振動は場の量子である。その量子が不安定なら、光子、電子、陽電子、その他の身近な粒子に崩壊しうる。}
\end{classroomqa}

\section{スローロールとハッブル摩擦}

場が縁を越えてしまったら、なぜ永遠に揺れ続けないのか。なぜ前後に揺れ、場合によっては再び登り返さないのか。ここで与えられる答えは宇宙論的減衰機構である。それは真の微視的摩擦ではないが、運動方程式の中ではそのように振る舞う。

空間勾配を無視し、場を一様とみなす。すると運動項とポテンシャル項は\(T-V\)の力学を持つ単一の自由度のように見えるが、それらはエネルギー密度であるため、物理的体積\(a^3\)を掛けなければならない。その結果、次の形の一様有効ラグランジアンが導かれる。
\begin{equation}
L_{\mathrm{hom}}=a^3(t)\left[\frac12\dot\phi^2-V(\phi)\right].
\end{equation}
\editorialnote{黒板の一番上の行は一部が切れている。この一様ラグランジアンは、その下に書かれた確実な運動方程式を再現するのに必要な範囲でのみ再構成されている。}

黒板上で確実な方程式は
\begin{equation}
\frac{d}{dt}\left(a^3\dot\phi\right)=-a^3\frac{\partial V}{\partial\phi}.
\end{equation}
である。時間微分を展開すると
\begin{equation}
a^3\ddot\phi+3a^2\dot a\,\dot\phi=-a^3\frac{\partial V}{\partial\phi},
\end{equation}
となり、\(a^3\)で割ると
\begin{equation}
\ddot\phi+3H\dot\phi=-\frac{\partial V}{\partial\phi}.
\end{equation}
が得られる。マイナス符号は口頭の導出で明示的に修正されている。力はポテンシャルエネルギーの導関数にマイナスを付けたものである。

\begin{figure}[t]
\centering
\includegraphics[width=0.78\textwidth]{lecture_08_figure_05.png}
\caption{ハッブル摩擦項\(3H\dot\phi\)を示す一様インフラトン運動方程式。}
\lectureframe{00:57:10}
\end{figure}

\(3H\dot\phi\)という項がハッブル摩擦と呼ばれるものである。それは完全に因子\(a^3\)を微分することから生じる。しかし力学的には、水中を落下する岩の運動における抗力項のように見える。力の項、加速度の項、そして速度に線形な項がある。場はすぐに終端速度に達し、その後はプラトーに沿ってゆっくりと動く。これがスローロール領域である。傾きが浅く場がゆっくり動く間、ポテンシャルエネルギーはゆっくりとしか変化しないため、膨張はほぼ指数関数的にとどまる。

最小値付近では場は振動する。それらの振動がインフラトン量子である。量子が不安定であれば崩壊し、そのようにしてプラトー上のエネルギーが後の宇宙の高温物質と放射の内容に変わる。

\begin{classroomqa}
\lecturetimestamp{00:50:48}
\audiencequestion{なぜ場は永遠に振動し続けないのか、そして何が減衰を与えるのか。}
\lecturerresponse{膨張そのものが摩擦のような項を与える。一様ラグランジアン中の因子\(a^3\)が\(3H\dot\phi\)項を生み、それが粘性抗力のように振る舞う。}
\end{classroomqa}

\section{量子揺らぎとホライズン退出}

ここで明らかな問題が生じる。もしインフレーションが宇宙を極めて滑らかにするなら、なぜ銀河が必要とするまさにその種を消し去らないのか。答えは量子力学である。生き残る揺らぎは不完全なアイロンがけの名残ではない。それらは新たに生成される。

しばらく膨張を忘れよう。たとえ場がほぼ完全に一様であっても、量子力学は\(\phi\)と\(\dot\phi\)の両方を正確に固定することを許さない。基底状態もなお揺らいでいる。これは通常の調和振動子が零点エネルギーを持つのと同じ理由である。日常的な状況では、それらの揺らぎは短波長で速く振動するため見えにくく、観測はそれらを平均してしまう。

膨張は話を変える。共動座標では場のモードはほぼ同じ座標波長を保つが、物理的波長はスケール因子とともに成長する。
\begin{equation}
\lambda_{\rm phys}=a(t)\,\Delta x.
\end{equation}
インフレーション中、物理的ハッブルスケールはほぼ固定されており、
\begin{equation}
d_H\sim H^{-1},
\end{equation}
したがって対応する共動ホライズンは次のように縮む。
\begin{equation}
d_{H,\mathrm{coord}}\sim \frac{1}{aH}.
\end{equation}

ロープの類推は凍結を直感的にする。波が振動するのは、ロープの各部分が隣接する部分に引っ張られるからである。しかし波長がホライズンと同じくらい大きくなると、関連する隣接部分は因果的接触の外にある。復元力の引っ張りが消える。モードがホライズンより長くなると、それは凍結する。

これは繰り返し起こる。より長い波長が先に凍結する。より短い波長は後で凍結する。たとえ微小な波長で始まったモードでも、最終的には引き伸ばされてホライズンを横切る。したがってインフレーションは単に構造を消し去るのではない。凍結した場の変動を生み続けるのである。

この過程はインフレーションが長く続くからといって止まるわけでもない。古い凍結モードは巨大なスケールに引き伸ばされるが、新しいモードは常に凍結しつつある。蓄積されたパターンは、宇宙論に関連する意味で統計的に定常になる。これがほぼスケール不変なスペクトルの起源である。同じホライズン横断過程が連続する対数スケールで繰り返され、ある波長帯から次の波長帯へとほぼ同じ振幅を生み出すのである。

場がプラトーを転がり落ちた後、地平線は再び成長する。一時的なインフレーション期の地平線の外にあったモードは、後に再び中へ入ってくる。通常の言葉では、それらはインフレーション中に地平線を出て、その後再び入る。その頃にはそれらは巨視的なスケールに引き伸ばされているので、再び振動し始めるときには、非常に大きな構造として振動する。

\begin{classroomqa}
\lecturetimestamp{01:12:02}
\audiencequestion{時空グラフ上で地平線はどのように見えるか。}
\lecturerresponse{指数関数的膨張の間、地平線までの物理的距離はほぼ\(c/H\)に固定されたままであるが、共動座標のグラフ上では、座標スケールが引き伸ばされているため、その固定された物理的距離はますます小さく見える。共動座標で言えば、地平線は\(1/(aH)\)のように縮む。}
\end{classroomqa}

\begin{classroomqa}
\lecturetimestamp{01:15:59}
\audiencequestion{揺らぎが引き伸ばされるとき、振幅はどうなるのか、また定常状態はあるのか。}
\lecturerresponse{個々のモードは引き伸ばされるにつれて希薄になるが、新しいモードが次々と凍結し続ける。その結果は、揺らぎが完全に消えることではなく、ほぼスケール不変な統計的定常パターンである。}
\end{classroomqa}

\section{凍結した場の変動から密度変動へ}

ひとたび場が場所ごとに小さな変動を獲得すると、異なる領域がわずかに異なる時刻に縁を越えるとき、それらの変動は密度差に変換される。少しだけ坂を下った領域が最初に崖に到達する。その真空エネルギーが放射に変換されれば、その放射は膨張とともに希薄化する。まだプラトー上にある隣の領域は真空優勢のままであり、希薄化しない。二番目の領域が落ちる頃には、最初の領域はすでにそのエネルギーの一部を赤方偏移させて失っている。

だからこそ、より遅い転がりはより大きな最終密度コントラストを与える。場がプラトーに沿って非常にゆっくりとしか動かなければ、場の値の小さな差は、ある領域と次の領域が崖に到達する間のより大きな時間差に対応する。先に落ちる領域は、後に落ちる領域がそのエネルギーを変換する前に、より多くの時間を希薄化に使える。傾きが平坦であるほど終端速度は小さくなり、最終密度変動は大きくなる。より高いプラトーもまた、より大きな揺らぎを与える傾向がある。プラトーの幅は主に、インフレーションがどれだけ続くかを制御する。

揺らぎのデータによって厳しく制約されるのは、高さと傾きの組み合わせだけである。エネルギー・スケールが十分高ければ、目に見える重力波の寄与もあるはずである。数値的な限界は、おおよその目安としてのみ与えられている。おおよそ言えば、プラトー上のエネルギー密度がプランク単位でおよそ\(10^{-14}\)よりずっと大きければ、宇宙マイクロ波背景放射はすでに過剰な重力波構造を示しているはずであり、およそ\(10^{-16}\)よりずっと小さければ、そのような信号を見る現実的な見込みはほとんどないであろう。\editorialnote{話し手はこれらの数値を明示的に近似として扱い、正確な値を覚えていないと述べている。それらは桁の程度の記述としてのみここに保持されている。}

単位は通常の自然単位の流儀で説明される。質量\(m\)の粒子のコンプトン波長は
\begin{equation}
\lambda_C=\frac{\hbar}{mc}.
\end{equation}
である。\(c=\hbar=1\)ならば、質量と逆長さは同じ単位を持つ。エネルギー密度は体積あたりのエネルギーであるから、質量の四乗の単位を持つ：
\begin{equation}
[\rho]=(\mathrm{mass})^4.
\end{equation}

一点で見れば、凍結したモードの蓄積はランダムウォークのように見える。空間の異なる点は異なるランダムウォークを経験し、まさにそれによって量子ノイズが空間的変動に変わる。

\begin{classroomqa}
\lecturetimestamp{01:33:04}
\audiencequestion{\(\phi\)はランダムウォークを経験するのか、そしてそれによって以前の出発点は無関係になるのか。}
\lecturerresponse{任意の固定点では、連続する凍結モードがランダムな位相で加わるので、\(\phi\)は確かにランダムウォークを行う。そのまさに成功が困難でもある。インフレーションは以前の情報を後の揺らぎで置き換え続けるので、インフレーション以前の段階を回復するのは難しくなる。}
\end{classroomqa}

\section{インフレーションが隠すもの}

この理論は非常に特定の意味で成功している。それは広大なスケール範囲にわたる揺らぎを一つの根底にある機構と結びつける。しかし同じ機構が、それ以前にあったものを洗い流してしまう。以前の不規則性は引き伸ばされて消え、後に生成された量子揺らぎが関連を持つものとして残る。だからこそインフレーションは魅力的であり、同時に歯がゆい。それは滑らかさと揺らぎのスペクトルを説明するが、以前の歴史の多くも消し去ってしまう。

このことは出発点についての不安を引き起こす。非常に高いエネルギー密度が問題なのではない。問題は、プラトーのエネルギーが自然にプランク的でも自然に微小でもないように見えることである。通常の単位では高いが、プランク単位では依然として小さい。基礎理論の明白な最初の停止点のようには見えない。

また、この単純な描像には文字通りの\(t=0\)は存在しない。インフレーションが長く続きすぎると、それ以前の段階に関するほとんどすべての情報は失われる。e-foldingsの数が必要最小限をわずかに上回る程度であれば、前の段階の痕跡がいくらか生き残るかもしれないが、一般にインフレーションは過去を消し去る。

同じ注意はド・ジッター空間の扱いにも現れる。純粋な指数法則は後期におけるインフレーション分枝に対して正しいが、厳密なド・ジッター解は次でよりよく記述される。
\begin{equation}
a(t)\propto \cosh(Ht).
\end{equation}
この解は最小サイズまで収縮し、その後膨張する。収縮分枝を宇宙論において文字通りに取るべきかどうかは未解決のままである。

\begin{classroomqa}
\lecturetimestamp{01:37:31}
\audiencequestion{指数関数的振る舞いが過去へ無限に続き、始まりがなかったということはあり得たのか。}
\lecturerresponse{この単純な描像ではあり得ない。厳密なド・ジッター空間は全時間にわたって永遠に膨張する指数関数ではなく、収縮分枝と膨張分枝をもつ\(\cosh(Ht)\)のようなものである。その初期の分枝を文字通りに取るべきかどうかは未解決の問題として残される。}
\end{classroomqa}

\section{調整、曲率、人間原理的限界}

第一近似では、重要なパラメータはプラトーの高さ、その傾き、そしてその幅\(\Delta\phi\)である。傾きは終端速度を定め、したがって揺らぎの大きさを定めるため重要である。幅はインフレーションがどれだけ続くかの決定に寄与するため重要である。基本的な場の変位はプランク質量よりずっと大きくなるべきではない、またエネルギー密度はプランク単位で1程度を超えるべきではない、というのが一般的な先入観である。そのような粗い限界の下では、十分なインフレーションはパラメータ空間全体を埋めない。ここで引用される見積もりは、傾きをおよそ\(60\)から\(62\)のe-foldingsに対して十分浅くするための、数パーセント、およそ\(2\%\)の調整である。

\begin{classroomqa}
\lecturetimestamp{01:46:02}
\audiencequestion{インフレーションはパラメータ空間の観点で理解できるのか、またどれほど調整されているのか。}
\lecturerresponse{はい。第一近似では、プラトーの高さ、傾き、幅の観点で考える。講義は、十分なe-foldingsを得るには数パーセント、およそ\(2\%\)の水準の調整が必要であるという粗い見積もりを与えている。}
\end{classroomqa}

その軽度の調整は、次に、桁で\(10^{-123}\)の現在の宇宙定数というはるかに鋭い難問と比較される。議論は人間原理的問題へと向かう。どちらの問題も、構造が形成されるために何が必要かを問うことで言い換えられるからである。

正の宇宙定数は、ニュートン的描像では距離に比例する普遍的な斥力のように働く。それがずっと大きければ、今日銀河を吹き飛ばすことまでしなくても致命的であり得る。原始の\(10^{-5}\)揺らぎが銀河へと凝縮するのを妨げるのに十分大きければよいだけである。挙げられている例は、観測値よりおよそ2桁大きい値で既にそうなるというものである。これがワインバーグの議論に関連する点である。宇宙定数は正確にゼロである必要はなく、構造を許すのに十分小さいだけでよい。

負の曲率も同じ方向に働く。以前の講義のニュートン的言語では、\(k<0\)は脱出速度より速い後退を意味する。外向きの推力が大きすぎても、小さな原始の過密度が銀河へと変わるのを妨げる。インフレーションはその曲率を薄める。引用されている下限はおよそ\(59.5\)のe-foldingsである。次に\(62\)程度の値が、しきい値を安全に上回る代表的なインフレーション史として取られる。

これは確率の問いを変える。十分なインフレーションを与えるパラメータ空間の生の体積を問うなら、答えは数パーセントにすぎないかもしれない。銀河が形成され、問いを発する観測者が存在する宇宙に条件付けるなら、条件付き確率はずっと大きくなり得る。そこで議論は閉じられる。最終的な解決ではなく、微調整の難問、それらの人間原理的再解釈、そしてここに大きな論争が残っているという認識をもって。